\section{The chain constraint: when forest optimization disappears}
\label{sec:chain}

The two-graph model has an additional constraint on native geometric inputs.
Its cap incidences and raw attaching counts are parts of the same cellular
chain complex. This constraint makes a substantial part of the apparent
forest-choice problem collapse to one tree extension. It also provides a
necessary realizability check for proposed abstract examples.

\subsection{Raw parity and complementary tree coordinates}

Fix an authenticated core $K$ and a geometric forest contraction among its
zero- and one-handles. Let $E$ index its $m$ two-handles, let $C$ index the
remaining one-handles, and let $q$ be the number of three-handles. Write
$a_{ec}$ for the \emph{uncapped} geometric occurrence count of one-handle
$c$ in the attaching cycle of two-handle $e$. Define
\[
 A=(a_{ec}\bmod 2)\in\F^{E\times C},\qquad
 B\in\F^{E\times\{1,\ldots,q\}},
\]
where $B_{ev}$ counts the cap occurrences of $e$ at three-handle $v$ modulo
two. Both ends of a loop are counted. Orientations disappear over $\F$,
but geometric occurrences must be counted before reduction modulo two.

The handle chain complex gives
\begin{equation}
 \partial_2=A^{\mathsf T},\qquad
 \partial_3=B,\qquad A^{\mathsf T}B=0.
 \label{eq:boundary-boundary}
\end{equation}
The matrix $B^{\mathsf T}$ is the reduced incidence matrix of the cap graph
$G$, with the formal exterior root omitted. By \cref{thm:top}, $G$ is
connected in the native source setting. Hence $\rk B=q$, and
\begin{equation}
 \beta_2(K;\F):=\dim_{\F}H_2(K;\F)=m-\rk A-q.
 \label{eq:chain-beta}
\end{equation}
Several core components cause no problem when each reaches an exterior
end: the common graph root is formal and the three-handle columns remain
independent. Closed components with nonzero $H_3$ require a different rank
statement.

Suppose first that every \emph{raw} row has length at most two. Construct
the occurrence graph $H$ as in \cref{sec:joint}: singleton rows join their
column to a ground vertex, two distinct occurrences join their columns,
and double occurrences or empty rows are loops. Then $A^{\mathsf T}$ is
exactly a reduced incidence representation of $H$. Write $r_H(S)$ for its
graphic rank on edge set $S$, and $\nu_H(S)=|S|-r_H(S)$ for its cycle
nullity.

\begin{theorem}[Geometric short-row invariance]
\label{thm:chain-invariance}
If $\sum_c a_{ec}\leq2$ for every $e\in E$, then every spanning tree $F$ of
the cap graph satisfies
\begin{equation}
 r_H(E\setminus F)=r_H(E),\qquad
 \nu_H(E\setminus F)=\beta_2(K;\F).
 \label{eq:tree-invariance}
\end{equation}
Consequently, one such residue admits complete unit-column peeling if and
only if every such residue does. Changing the top-handle tree cannot
improve its occurrence-graph nullity.
\end{theorem}

\begin{proof}
Transposing \cref{eq:boundary-boundary} gives
\[
 \operatorname{col}(A)\subseteq\ker(B^{\mathsf T}),
\]
the cycle space of $G$. Projection onto the coordinates $E\setminus F$ is
injective on this cycle space: a vector in its kernel would be a cycle
supported entirely on a forest, and therefore zero. The same projection
is injective on $\operatorname{col}(A)$, so $\rk A_{E\setminus F}=\rk A$.
These ranks are respectively $r_H(E\setminus F)$ and $r_H(E)$.

Every cap spanning tree has $q$ edges. Substitution in
\cref{eq:chain-beta} proves the nullity identity. Finally, the short-row
peeling criterion of \cref{sec:joint} identifies complete peeling with
zero occurrence-graph nullity.
\end{proof}

This is a restriction on the optimization problem, rather than a defect
in the abstract matroid theorem. In compatible geometric inputs, the row
space representing $H$ lies in the cycle space of $G$. The two graphs
cannot be varied independently.

\subsection{Mandatory rows and the homological criterion}

Capping counts at two preserves singleton predicates but can change
parity. For example, a raw row $(3)$ has binary row $(1)$, whereas its
capped row $(2)$ represents a loop and hence a zero incidence column.
With no three-handles, this example has $m-\rk A-q=0$ but occurrence-graph
nullity one. Thus \cref{thm:chain-invariance} cannot simply be applied to
arbitrary capped short rows.

There is an exact repair for the restricted strategy. Put
\begin{align}
 D_{\mathrm{long}}&=\bigl\{e:\sum_c\min(2,a_{ec})>2\bigr\},\notag\\
 D_{\mathrm{loop}}&=\{e:a_e=0\text{ or }a_e=k\mathbf e_c
                         \text{ for some }k\geq2\},\notag\\
 \overline D&=D_{\mathrm{long}}\cup D_{\mathrm{loop}}.
 \label{eq:mandatory-rows}
\end{align}
The first set is required to be removed by the restricted rule in
\cref{sec:joint}. The second consists of all occurrence-graph loops, none
of which can remain in a completely peelable residue. Thus every successful
certificate in that strategy must place all of $\overline D$ in its cap
tree. Every remaining row is genuinely one unit occurrence or two distinct
unit occurrences; its raw and capped parities agree.

\begin{theorem}[Mandatory-forest reduction]
\label{thm:mandatory-chain}
Assume $\overline D$ is a forest of $G$. Set $k=|\overline D|$,
$E_1=E\setminus\overline D$, and let $H_1$ be the occurrence graph on $E_1$.
For every cap spanning tree $F\supseteq\overline D$,
\begin{equation}
 \nu_{H_1}(E\setminus F)
   =(m-k)-(q-k)-r_{H_1}(E_1)
   =\beta_2(K;\F).
 \label{eq:mandatory-beta}
\end{equation}
For the fixed primal forest and the restricted short-row rule, a complete
top-deletion and raw-peeling certificate exists exactly when
$\overline D$ is a cap forest and $\beta_2(K;\F)=0$.
\end{theorem}

\begin{proof}
Let $Q$ have one column for each component of $G[\overline D]$ not
containing the exterior root. Its entries are one on the three-handle
vertices in that component and zero elsewhere. Thus
\[
 B_{\overline D}Q=0,\qquad B_1:=B_{E_1}Q
\]
is the transpose of the reduced incidence matrix of $G/\overline D$.
Multiplying \cref{eq:boundary-boundary} by $Q$ gives
\begin{equation}
 A_{E_1}^{\mathsf T}B_1=0.
 \label{eq:contracted-chain}
\end{equation}
Because $\overline D$ is a forest, the contracted connected graph has
$q-k+1$ vertices. The projection argument in
\cref{thm:chain-invariance} applies to $A_{E_1},B_1$. It shows that
restriction away from the tree $F\setminus\overline D$ preserves
$r_{H_1}(E_1)$, proving the first equality in \cref{eq:mandatory-beta}.

To identify this value with the original core's homology, use that
$B_{\overline D}$ has full row rank $k$. Choose a right inverse $R$ with
$B_{\overline D}R=I_k$. The original chain identity implies
\[
 A_{\overline D}^{\mathsf T}
     =A_{E_1}^{\mathsf T}B_{E_1}R.
\]
The deleted column maps of $\partial_2$ are therefore combinations of the
surviving ones. Consequently
\[
 \rk A=\rk A_{E_1}=r_{H_1}(E_1).
\]
Equation~\eqref{eq:chain-beta} now proves the second equality.

Every successful restricted certificate has a cap tree containing
$\overline D$, so this set must be a forest. Its residue has nullity zero,
forcing $\beta_2=0$. Conversely, if both conditions hold, extend
$\overline D$ to any cap spanning tree. Its residue is an occurrence
forest, and therefore peels completely. The rooted schedule of
\cref{lem:rooted} and the authenticated geometric cancellation rules turn
these graph steps into a core cancellation certificate.
\end{proof}

The contraction in this proof chooses an edge set; it does not assert that
the mandatory edges alone form an initial legal cancellation sequence.
After extension to $F$, the complete geometric trace is ordered from the
exterior root as in \cref{lem:rooted}.

\begin{corollary}[One extension suffices]
\label{cor:single-extension}
On supplied authenticated incidence data, the restricted strategy can be
decided by testing whether $\overline D$ is a cap forest, extending it once
to a spanning tree when possible, and testing or peeling the residual
occurrence graph. No search over alternative extensions is required.
\end{corollary}

If the incidence lists are indexed and isolated occurrence vertices are
suppressed, let $N$ be the number of nonzero raw incidence entries. Finding
$\overline D$, traversing its components, building the quotient, choosing
a spanning tree and peeling use $O(m+q+N+1)$ graph work and memory.
Component traversals suffice; a general binary-matrix rank computation is
unnecessary. Source authentication, index construction and the bit costs
of reading or comparing binary counts are separate from this graph bound.
The positive output remains an ordinary replayable cancellation trace.

This is not a generic handlebody test based on homology. Sufficiency uses
the authenticated handle charts, the fixed primal forest, the mandatory
forest hypothesis and the particular short-row rule. A vanishing second
Betti number without those hypotheses does not supply the cancellations.
Moreover, failure of the restricted strategy can occur for a handlebody.
The later one/two cancellations also need their own marking analysis if a
consumer requires arbitrary boundary-pattern transport.

\subsection{Consequences for abstract evidence and deficiency}

The retained abstract exchange regression makes the distinction concrete.
Its cap edges are
$e_0=v_0v_1$, $e_1=v_1\rho$, $e_2=v_0\rho$, and its raw rows on one
occurrence column are $(1),(0),(2)$. The cap tree $\{e_0,e_1\}$ leaves an
occurrence loop, whereas $\{e_1,e_2\}$ leaves a peelable singleton.
However, its binary matrices satisfy
\[
 B=\begin{pmatrix}1&1\\0&1\\1&0\end{pmatrix},\qquad
 A=\begin{pmatrix}1\\0\\0\end{pmatrix},\qquad
 A^{\mathsf T}B=(1,1)\ne0.
\]
Thus these exact data cannot be the incidences of an authenticated handle
core. They remain a useful test of the abstract exchange algorithm, but
provide no example of a native geometric speedup.

\begin{example}[A geometric family with several top trees]
\label{ex:geometric-cycle}
Start with a standard solid torus $V$, presented by one zero-handle and one
one-handle $c$. Attach $r\geq2$ two-handles along disjoint parallel primitive
\emph{longitudes} of $V$, each crossing the cocore of $c$ once. These curves
represent a generator of $\pi_1(V)$; they are not meridians of $V$.

For an explicit embedding, use the standard genus-one splitting
$S^3=V\cup W$, in which a meridian of the complementary solid torus $W$ is
a longitude of $V$. Take the two-handles to be small neighbourhoods of $r$
disjoint meridian disks in $W$. Their complement in $W$ consists of $r$
balls. The resulting core is therefore $S^3$ with $r$ open balls removed,
equivalently a ball with $r-1$ interior balls removed. Fill any $0\leq q\leq r-1$
of its spherical boundary components with three-handles.

The $r$ complementary chambers lie cyclically around $W$. Consequently
the cap graph is the $r$-cycle with the unfilled chamber vertices
identified with $\rho$; for $r=2$ the cycle has two parallel edges.
The occurrence graph has $r$ parallel edges between $c$ and its ground,
and rank one. Every cap tree deletes $q$ rows and hence leaves
\[
 |E\setminus F|=r-q,\qquad
 \nu_H(E\setminus F)=r-q-1=\beta_2(K;\F).
\]
When $q=r-1$, there are $r$ different cap trees and every one succeeds;
the filled core is a ball. When $r=3$ and $q=1$, there are two cap trees,
and each leaves nullity one. Thus genuine geometric inputs can have
several tree choices and positive invariant defect. The ports in this
family are empty; the construction supplies geometric examples only,
without a native triangulation benchmark claim.
\end{example}

The exact abstract deficiency of \cref{sec:joint} remains meaningful.
For any cap spanning tree $F$, the rows $B_F$ form an invertible square
matrix. Equation~\eqref{eq:boundary-boundary} then expresses $A_F^{\mathsf T}$
in terms of $A_{E\setminus F}^{\mathsf T}$, so every such restriction retains
the raw binary rank $\rk A$. When only $D_{\mathrm{long}}$ is forced,
passing from the remaining raw short rows to their capped occurrence graph
merely zeros the columns arising from odd single multiplicities at least
three. It cannot increase rank. Every corresponding occurrence nullity is
therefore at least $\beta_2(K;\F)$.

If $\overline D$ is a forest, \cref{thm:mandatory-chain} supplies an extension
attaining this lower bound. In that case the abstract minimum deficiency
equals the second Betti number as well. If $\overline D$ is cyclic, a
complete short-row certificate is impossible; optimizing a partial
residue can still be a diagnostic. Further native progress must then use
information outside this complete short-row forest strategy, such as the
structure of retained longer rows, a different primal contraction, or
authenticated geometric moves that change the occurrence data.
